% Quadro de sinais alternados de cofatores em matriz 3x3.
% Usado em: capitulo_01_calculo-de-determinantes.md (seção 4.2)
\begin{tikzpicture}[
  cell/.style={
    rectangle,
    minimum width=12mm,
    minimum height=12mm,
    inner sep=0pt,
    draw=gray!60,
    thick
  },
  positivo/.style={cell,fill=eleveBlue!12},
  negativo/.style={cell,fill=eleveAccent!12}
]
  % Linha 1
  \node[positivo] (n11) at (0,2)   {$+$};
  \node[negativo] (n12) at (1.2,2) {$-$};
  \node[positivo] (n13) at (2.4,2) {$+$};

  % Linha 2
  \node[negativo] (n21) at (0,0.8)   {$-$};
  \node[positivo] (n22) at (1.2,0.8) {$+$};
  \node[negativo] (n23) at (2.4,0.8) {$-$};

  % Linha 3
  \node[positivo] (n31) at (0,-0.4)   {$+$};
  \node[negativo] (n32) at (1.2,-0.4) {$-$};
  \node[positivo] (n33) at (2.4,-0.4) {$+$};

  % Rótulo
  \node[font=\small,below=4mm of n32] {Sinais de $(-1)^{i+j}$};
\end{tikzpicture}
